\chapter{Annexes}

\section{Scripts SQL}
Extrait de requête de contrôle qualité exécutée sous pgAdmin pour vérifier la structure d'une table de faits via le catalogue système~:
\begin{verbatim}
SELECT column_name, data_type, is_nullable
FROM information_schema.columns
WHERE table_schema = 'DW' AND table_name = 'Fact_Comptage'
ORDER BY ordinal_position;
\end{verbatim}

Extrait de requête de contrôle de cohérence référentielle (détection de clés étrangères orphelines)~:
\begin{verbatim}
SELECT f.*
FROM "DW"."Fact_deplacement" f
LEFT JOIN "DW"."Dim_Individu" d ON f."FK_id_individu" = d."PK_id_individu"
WHERE d."PK_id_individu" IS NULL;
\end{verbatim}

\section{Jobs Talend}
Le tableau~\ref{tab:talend_jobs_annexe} récapitule l'ensemble des jobs Talend développés pour l'alimentation de l'entrepôt de données.

\begin{table}[H]
\centering
\caption{Jobs Talend développés}
\label{tab:talend_jobs_annexe}
\begin{tabular}{ll}
\toprule
\textbf{Job} & \textbf{Rôle} \\
\midrule
SA\_individu, SA\_menage, SA\_deplacement & Chargement de la Staging Area \\
Job\_dim\_site, Job\_dim\_geo & Construction des dimensions de référence \\
Job\_dim\_individu, Job\_dim\_menage & Construction des dimensions individu / ménage \\
Job\_deplacement & Construction de \texttt{Fact\_deplacement} \\
Job\_comptage & Construction de \texttt{Fact\_Comptage} \\
Job\_accessibilite & Construction de \texttt{fact\_access} \\
Job\_indivmen & Construction de \texttt{Fait\_IndividuMenage} \\
\bottomrule
\end{tabular}
\end{table}

\section{Captures Knowage}
Les sheets du cockpit Knowage BI (Accueil, Trafic, Déplacements, Démographie, Accessibilité, IA) sont illustrées par les figures du chapitre~\ref{ch:resultats} ainsi que par les captures du dossier \texttt{figures/kpis/} du présent rapport.

\section{Captures React}
Les pages de l'application web (espace citoyen et espace décideurs) sont illustrées par les figures du chapitre~\ref{ch:dev}, présentant respectivement le parcours citoyen (recommandation de mode de transport) et les tableaux de bord décideurs (zones à risque, anomalies, segmentation, simulateur de risque).

\section{Documentation de l'API}
Le tableau~\ref{tab:api_annexe} liste l'ensemble des points d'accès REST exposés par le backend FastAPI, avec leur méthode HTTP et leur statut de protection.

\begin{table}[H]
\centering
\caption{Liste complète des endpoints REST exposés par l'API}
\label{tab:api_annexe}
\begin{tabular}{llc}
\toprule
\textbf{Méthode} & \textbf{Endpoint} & \textbf{Protégé (JWT)} \\
\midrule
POST & /auth/login & -- \\
GET & / & Non \\
GET & /health & Non \\
GET & /quartiers & Non \\
POST & /recommander & Non \\
GET & /segments & Oui \\
GET & /modes & Oui \\
GET & /zones-risque & Oui \\
GET & /zones-risque/resume & Oui \\
GET & /api/segmentation/profils & Oui \\
GET & /api/satisfaction & Oui \\
POST & /api/feedback & Non \\
GET & /api/feedback/stats & Oui \\
GET & /api/anomalies/summary & Oui \\
GET & /api/anomalies/sites & Oui \\
GET & /api/anomalies/sites/\{site\_id\}/details & Oui \\
POST & /predict-inaccessibility & Oui \\
GET & /api/ml/features-importance & Oui \\
GET & /api/ml/metrics & Oui \\
\bottomrule
\end{tabular}
\end{table}

La documentation interactive complète (Swagger UI) est générée automatiquement par FastAPI et accessible à l'adresse \texttt{/docs} du serveur.

\section{Documentation des modèles de Machine Learning}
Chaque modèle est sérialisé avec ses métadonnées (liste des variables, métriques de performance, date d'entraînement) au format \texttt{.pkl} via \texttt{joblib}~: \texttt{kmeans\_segmentation.pkl}, \texttt{rf\_recommandation.pkl}, \texttt{inaccessibilite\_model.pkl}, ainsi que les fichiers de métadonnées associés (\texttt{metadata.json}, \texttt{zones\_risque.json}). Le détail méthodologique de chaque modèle (objectif, variables, algorithme, résultats, intégration) est présenté au chapitre~\ref{ch:ml}.

\section{Manuel utilisateur}

\subsection{Espace citoyen}
\begin{enumerate}
\item Accéder à la page d'accueil de l'application (\texttt{WelcomePage})~;
\item Cliquer sur «~Planifier mon trajet~» pour accéder au formulaire de recommandation~;
\item Renseigner son profil de mobilité (zone de résidence, âge, statut, disponibilité d'un véhicule, fréquence des déplacements)~;
\item Consulter le \emph{top 3} des modes de transport recommandés, avec durée et coût estimés.
\end{enumerate}

\subsection{Espace décideurs}
\begin{enumerate}
\item Accéder à la page de connexion (\texttt{LoginDecideurs}) et s'authentifier avec ses identifiants (chef d'exploitation ou directeur de planification)~;
\item Consulter la vue générale (\texttt{VueGenerale}) pour un aperçu synthétique des indicateurs clés~;
\item Naviguer vers les modules thématiques selon ses besoins~: zones à risque, anomalies de trafic, segmentation des usagers, simulateur de risque, cockpit Knowage~;
\item Pour le directeur de planification uniquement~: consulter la page de transparence des modèles (\texttt{MlInsights}) pour auditer les performances et l'importance des variables~;
\item Exporter un rapport PDF de synthèse depuis les pages \emph{Zones à risque} ou \emph{Anomalies} si nécessaire.
\end{enumerate}

\section{Dictionnaire des codes de variables EMD utilisés}
\begin{table}[H]
\centering
\caption{Principaux codes de variables EMD exploités}
\begin{tabular}{ll}
\toprule
\textbf{Code} & \textbf{Signification} \\
\midrule
I2, I4, I9 & Caractéristiques géographiques du domicile (quartier, strate, district) \\
M21 & Taille du ménage \\
M55 & Manque de transport en commun déclaré \\
M66 -- M70 & Accessibilité de base (distance, durée de marche, modes disponibles) \\
M68 & Exposition aux inondations \\
M69 & Résilience d'accès en cas de pluie \\
M73, M87 & Problèmes majeurs déclarés dans le quartier \\
M90, M91 & Accès aux écoles et daaras \\
\bottomrule
\end{tabular}
\end{table}

\section{Glossaire}
\begin{description}
\item[BI] Business Intelligence -- informatique décisionnelle.
\item[CETUD] Conseil Exécutif des Transports Urbains de Dakar.
\item[DBSCAN] Density-Based Spatial Clustering of Applications with Noise -- algorithme de clustering spatial.
\item[EMD] Enquête Ménages-Déplacements.
\item[ETL] Extract-Transform-Load -- processus d'extraction, transformation et chargement des données.
\item[GIMSI] Méthodologie de conception de tableaux de bord décisionnels.
\item[JWT] JSON Web Token -- jeton signé utilisé pour l'authentification.
\item[KPI] Key Performance Indicator -- indicateur clé de performance.
\item[LOF] Local Outlier Factor -- algorithme de détection d'anomalies par densité locale.
\item[RBAC] Role-Based Access Control -- contrôle d'accès basé sur les rôles.
\item[SHAP] SHapley Additive exPlanations -- méthode d'explicabilité des modèles de Machine Learning.
\item[SMOTE] Synthetic Minority Over-sampling Technique -- technique de rééquilibrage de classes.
\end{description}

\section{Liste des technologies utilisées}
\begin{table}[H]
\centering
\caption{Stack technique du projet}
\begin{tabular}{ll}
\toprule
\textbf{Couche} & \textbf{Technologie} \\
\midrule
Préparation des données & Python, Pandas, Google Colab \\
ETL & Talend Open Studio \\
Entrepôt de données & PostgreSQL \\
Restitution décisionnelle & Knowage BI \\
Machine Learning & scikit-learn, imbalanced-learn, SHAP, joblib \\
Backend applicatif & FastAPI, Pydantic, PyJWT, uvicorn \\
Frontend applicatif & React.js, Recharts, Leaflet, Framer Motion, jsPDF \\
Gestion de projet & Scrum \\
\bottomrule
\end{tabular}
\end{table}
